\section{The Bailout Protocol}\label{sec:bailout}

The Bailout Protocol is the mandatory exception-escalation mechanism of the \bxthree{} Framework. It is not an error-handling routine selected at runtime, nor an optional escalation pathway available to a system that chooses to invoke it. It is an architectural compulsion: any unresolved exception state encountered by any node in a \bxthree{} system \emph{must} propagate upward through the accountability chain until it reaches a named human principal, and no machine actor along that chain may absorb, resolve, or suppress the exception. This section provides the complete formal specification: the deterministic state machine governing exception propagation, the precise conditions governing each transition, the bail-out trigger condition, the escalation algorithm, the bypass mechanism enforcing the no-machine-actor property, and the timeout handling that guarantees liveness under all failure conditions.

The protocol operates against the backdrop of the \bxthree{} three-layer model. The \purpose{} layer carries intent, judgment, and final authority; it defines the operating range $R(n)$ for each node $n$ and receives all escalated exceptions. The \bounds{} layer carries bounded reasoning and constrained execution within $R(n)$; it evaluates proposed actions and initiates bailout when they exceed provisioned scope. The \fact{} layer carries deterministic enforcement and forensic auditability; it verifies proposed actions against the authorization ledger, enforces the state machine's transition relation, and maintains the append-only Forensic Ledger. Each spawned node carries an immutable parent pointer $\alpha(n)$, and every node carries a reference to its human accountability anchor $h(n)$ — a named human principal or institutional actor — established at node provisioning and immutable for the node's lifetime.

% --------------------------------------------------------------
\subsection{State Machine Formalism}\label{sec:state-machine}

The Bailout Protocol is formally modeled as a deterministic finite state machine
\[
\mathcal{B} = (Q,\ q_0,\ \Sigma,\ \rightarrow,\ Q_F)
\]
where:

\begin{itemize}\itemsep0em
\item $Q = \{\,\text{IDLE},\ \text{BOUNDED},\ \text{BAIL\_PENDING},\ \text{ESCALATING},\ \text{HUMAN\_ACKNOWLEDGED},\ \text{RESOLVED}\,\}$ is the finite set of protocol states;
\item $q_0 = \text{IDLE}$ is the designated initial state, entered at node startup and after every successful directive resolution;
\item $\Sigma$ is the input alphabet, comprising trigger events $e_{\text{trigger}}$, authorization confirmations $e_{\text{auth}}$, timeout signals $e_{\text{timeout}}$, human acknowledgments $e_{\text{ack}}$, binding resolution directives $e_{\text{resolve}}$, and rejection directives $e_{\text{reject}}$;
\item $\rightarrow \subseteq Q \times \Sigma \times Q$ is the deterministic transition relation (defined exhaustively in Table~\ref{tab:transition-relation});
\item $Q_F = \{\text{RESOLVED}\}$ is the set of terminal accepting states.
\end{itemize}

A node's state $q \in Q$ is stored in the node's \fact{} layer worksheet and is observable by its parent node and by the Bailout Monitor. State updates are atomic; no intermediate or partially-executing states are observable by any external actor. The machine is \emph{deterministic}: for every $(q, \sigma) \in Q \times \Sigma$, there exists at most one $q' \in Q$ such that $(q, \sigma, q') \in \rightarrow$. Determinism is enforced by the Bailout Monitor's priority function $\pi$, which totally orders simultaneous trigger conditions and selects the highest-priority enabled transition before committing any state change. If no transition is defined for a $(q, \sigma)$ pair, the machine stalls and records a protocol violation in the Forensic Ledger.

\begin{table}[htbp]
\caption{Bailout Protocol state transition relation $\rightarrow \subseteq Q \times \Sigma \times Q$. Rows are current state; columns are input events. Entries denote the next state. Blank cells indicate no defined transition.}
\label{tab:transition-relation}
\begin{tabular}{@{}lcccccc@{}}
\toprule
 & \multicolumn{6}{c}{Input event} \\
\cmidrule(l){2-7}
Current state & $e_{\text{trigger}}$ & $e_{\text{auth}}$ & $e_{\text{timeout}}$ & $e_{\text{ack}}$ & $e_{\text{resolve}}$ & $e_{\text{reject}}$ \\
\midrule
IDLE & BOUNDED & — & — & — & — & — \\
BOUNDED & — & HUMAN\_ACKNOWLEDGED & BAIL\_PENDING & — & — & — \\
BAIL\_PENDING & — & — & ESCALATING & — & — & — \\
ESCALATING & — & HUMAN\_ACKNOWLEDGED & ESCALATING & IDLE & RESOLVED & RESOLVED \\
HUMAN\_ACKNOWLEDGED & — & — & — & IDLE & RESOLVED & IDLE \\
RESOLVED & — & — & — & — & — & — \\
\bottomrule
\end{tabular}
\end{table}

\subparagraph{State definitions.}

\begin{itemize}\itemsep0em
\item[\textbf{IDLE.}] The node is operating within its provisioned authority. No exception condition is active. The \fact{} layer processes normal action requests through the standard dispatch path. The node may transition to BOUNDED at any time upon receipt of a trigger event satisfying the bail-out trigger condition (TC), defined in \S\ref{sec:bailout-trigger}.

\item[\textbf{BOUNDED.}] The node has detected an exception condition and is evaluating whether it is autonomously resolvable within current authority. The \bounds{} layer Bounds Engine is actively executing its resolution procedure. All machine-actor resolution paths remain open. If the Bounds Engine produces a resolution $a$ that the \fact{} layer approves before $T_{\text{bounds}}$ expires, the node returns to IDLE. If all resolution paths are exhausted or $T_{\text{bounds}}$ expires, the state advances to BAIL\_PENDING.

\item[\textbf{BAIL\_PENDING.}] The node has determined — either through explicit \texttt{bailout\_signal} from the Bounds Engine or through $T_{\text{bounds}}$ expiration without \fact{} layer-approved resolution — that the exception exceeds its current authority. The Bailout Protocol is formally activated, but propagation to $h(n)$ has not yet been initiated. This is the last architectural point at which a machine actor could theoretically suppress escalation. It cannot: entry to BAIL\_PENDING triggers an immediate atomic transition to ESCALATING, with no dwell time and no intermediate machine-actor intervention permitted.

\item[\textbf{ESCALATING.}] The exception is propagating upward through the parent chain toward $h(n)$. Each intermediate node receives the exception, appends its diagnostic context to the propagation chain, and forwards the exception to its parent via \texttt{SendToParent} without attempting local resolution. The state remains ESCALATING until either (a) the human anchor acknowledges receipt ($e_{\text{ack}}$), (b) the human anchor issues a binding directive ($e_{\text{resolve}}$ or $e_{\text{reject}}$), or (c) all pathways timeout after $N_{\max}$ retries at $T_{\text{cap}}$, propagating a \texttt{parent\_unreachable} event to $h(n)$ via an alternate functional pathway.

\item[\textbf{HUMAN\_ACKNOWLEDGED.}] The human anchor has received the exception, observed the diagnostic context, and acknowledged it. The node awaits a directive. No machine actor may issue, simulate, or substitute a directive in this state. The node remains in this state until the human principal issues \texttt{ACK} (return to IDLE) or a binding directive \texttt{RESOLVE} / \texttt{REJECT} (enter RESOLVED).

\item[\textbf{RESOLVED.}] The human principal has issued a directive and the protocol has recorded it in the authorization ledger. The directive is one of: \texttt{ACK} (continue with current parameters), \texttt{RESOLVE} (execute a specific binding resolution), or \texttt{REJECT} (disallow the proposed action and enter a quiescent error state). This state is terminal for the current protocol run.
\end{itemize}

\subparagraph{State invariants.} The following invariants hold for all reachable states of $\mathcal{B}$:

\begin{align}
\text{INV}_1 &: \forall n \in \text{nodes},\ \text{state}(n) \in Q \setminus \{\text{ESCALATING}\} \implies \neg\text{active\_bailout}(n). \tag{INV-1}
\end{align}
A node occupies a non-escalating state \emph{iff} it has no active bailout propagation in flight.

\begin{align}
\text{INV}_2 &: \text{state}(n) = \text{ESCALATING} \implies \text{active\_bailout}(n) \land \text{bailout\_target}(n) = h(n). \tag{INV-2}
\end{align}
When in ESCALATING, the bailout target is always and only the human anchor — no intermediate machine actor may be the terminus.

\begin{align}
\text{INV}_3 &: \text{state}(n) = \text{RESOLVED} \implies \exists d \in \text{directives}: d = \texttt{RESOLVE} \lor d = \texttt{REJECT} \lor d = \texttt{ACK}. \tag{INV-3}
\end{align}
Entry to RESOLVED always records a directive from $h(n)$, never from a machine actor.

% --------------------------------------------------------------
\subsection{Transition Conditions}\label{sec:transition-conditions}

Each transition in $\mathcal{B}$ is governed by a formal condition that must hold for the transition to fire. These conditions are enforced by the Bailout Monitor's pre-commit hook before any state change is recorded in the Forensic Ledger.

\begin{enumerate}
\item[\textbf{T1 — IDLE $\rightarrow$ BOUNDED (Trigger):}] The node observes an input $x$ satisfying the primary trigger condition TC:
\begin{equation}
\text{T1 fires} \iff \text{TRIGGER}(n, x) \land \text{state}(n) = \text{IDLE}. \tag{T1}
\end{equation}
Equivalently, an explicit \texttt{bailout\_signal} may fire T1 regardless of $x$:
\begin{equation}
\texttt{bailout\_signal}(x) \implies \text{T1 fires for node } n. \tag{T1-alt}
\end{equation}

\item[\textbf{T2 — BOUNDED $\rightarrow$ IDLE (Autonomous Resolution):}] The Bounds Engine produces a resolution $a$ that the \fact{} layer approves before $T_{\text{bounds}}$ expires:
\begin{equation}
\text{T2 fires} \iff \text{approved}(a) \land \text{state}(n) = \text{BOUNDED} \land t < T_{\text{bounds}}. \tag{T2}
\end{equation}
The \fact{} layer's \texttt{approved} predicate verifies that $a$ is consistent with the authorization ledger, that $a$ does not exceed $R(n)$, and that $a$ has not been previously rejected for this exception.

\item[\textbf{T3 — BOUNDED $\rightarrow$ BAIL\_PENDING (Timeout/Exhaustion):}] Either $T_{\text{bounds}}$ expires without \fact{} layer-approved resolution, or the Bounds Engine emits \texttt{bailout\_signal}:
\begin{equation}
\text{T3 fires} \iff \bigl(\text{state}(n) = \text{BOUNDED} \land t \geq T_{\text{bounds}} \land \neg\exists a: \text{approved}(a)\bigr) \lor \texttt{bailout\_signal}(x). \tag{T3}
\end{equation}

\item[\textbf{T4 — BAIL\_PENDING $\rightarrow$ ESCALATING (Protocol Activation):}] Immediately upon entry to BAIL\_PENDING, T4 fires atomically with no dwell time:
\begin{equation}
\text{T4 fires} \iff \text{state}(n) = \text{BAIL\_PENDING}. \tag{T4}
\end{equation}
There is no condition under which T4 may be suppressed or delayed by any machine actor.

\item[\textbf{T5 — ESCALATING $\rightarrow$ HUMAN\_ACKNOWLEDGED (Direct Human Contact):}] The Bailout Monitor receives an acknowledgment $e_{\text{ack}}$ from $h(n)$ via any functional pathway:
\begin{equation}
\text{T5 fires} \iff e_{\text{ack}} \in \Sigma \land \text{state}(n) = \text{ESCALATING} \land \text{sender}(e_{\text{ack}}) = h(n). \tag{T5}
\end{equation}

\item[\textbf{T6 — ESCALATING $\rightarrow$ ESCALATING (Pathway Timeout/Retry):}] A per-hop propagation timeout $T_{\text{prop}}$ expires without acknowledgment from the parent, and $N_{\max}$ retries have not been exhausted:
\begin{equation}
\text{T6 fires} \iff e_{\text{timeout}} \in \Sigma \land \text{state}(n) = \text{ESCALATING} \land \text{retries} < N_{\max}. \tag{T6}
\end{equation}
This transitions the node to ESCALATING with an updated parent pointer (advanced one hop toward $h(n)$) and a doubled propagation timeout per the backoff schedule.

\item[\textbf{T7 — ESCALATING $\rightarrow$ RESOLVED (All Pathways Exhausted):}] All $N_{\max}$ retries at $T_{\text{cap}}$ have been exhausted without acknowledgment, and the \texttt{parent\_unreachable} event has propagated to $h(n)$ via an alternate functional pathway:
\begin{equation}
\text{T7 fires} \iff \text{retries} \geq N_{\max} \land \text{state}(n) = \text{ESCALATING}. \tag{T7}
\end{equation}

\item[\textbf{T8 — HUMAN\_ACKNOWLEDGED $\rightarrow$ IDLE (Human Acknowledgment without Resolution):}] The human anchor issues \texttt{ACK} without a binding directive:
\begin{equation}
\text{T8 fires} \iff \text{directive} = \texttt{ACK} \land \text{state}(n) = \text{HUMAN\_ACKNOWLEDGED}. \tag{T8}
\end{equation}

\item[\textbf{T9 — HUMAN\_ACKNOWLEDGED $\rightarrow$ RESOLVED (Binding Directive):}] The human principal issues a \texttt{RESOLVE} or \texttt{REJECT} directive:
\begin{equation}
\text{T9 fires} \iff \text{directive} \in \{\texttt{RESOLVE}, \texttt{REJECT}\} \land \text{state}(n) = \text{HUMAN\_ACKNOWLEDGED}. \tag{T9}
\end{equation}
\end{enumerate}

% --------------------------------------------------------------
\subsection{Bail-Out Trigger Condition}\label{sec:bailout-trigger}

The bail-out trigger is the condition that causes a node to enter BOUNDED state (T1). Rather than requiring system designers to enumerate specific error codes in advance — a practice that provides coverage only for anticipated failure modes and leaves unanticipated ones invisible — the Bailout Protocol defines the trigger in terms of \emph{observational boundary violation}: a node escalates precisely when its observable input space contains a condition outside its provisioned operating range $R(n)$, and the node cannot produce a \fact{} layer-approved resolution within $T_{\text{bounds}}$.

\subparagraph{Primary trigger condition.} The trigger $\text{TRIGGER}(n, x)$ is defined as a conjunction of three sub-conditions, all of which must hold simultaneously:
\begin{equation}
\text{TRIGGER}(n, x) \;\equiv\; \underbrace{x \notin R(n)}_{(i)} \;\land\; \underbrace{\neg\text{resolved}(n, x)}_{(ii)} \;\land\; \underbrace{\text{confidence}(n, x) < \tau}_{(iii)} \tag{TC}\label{eq:trigger-condition}
\end{equation}
where:

\begin{itemize}\itemsep0em
\item[(i)] $x$ is an element of the node's observable input space that falls outside the provisioned operating range $R(n)$. This is the \emph{bound signal} condition: the node observes something it was not provisioned to handle autonomously.
\item[(ii)] $x$ has not been previously resolved by the node within the current operational session. This prevents re-trigger on known-exception inputs that have already received human adjudication.
\item[(iii)] $\text{confidence}(n, x) < \tau$ is the node's self-assessed confidence that it can produce a correct resolution for $x$ without human input, where $\tau$ is the provisioned confidence threshold. The confidence function $\text{confidence} : \text{node} \times \text{input} \rightarrow [0, 1]$ is defined at node provisioning and is \emph{not} adjustable at runtime by any machine actor. This prevents a machine actor from inflating its own confidence to suppress a warranted escalation.
\end{itemize}

\subparagraph{Trigger priority function.} When multiple trigger conditions fire simultaneously for node $n$, they are totally ordered by the priority function $\pi(n, x)$:
\begin{equation}
\pi(n, x) \;=\; w_1 \cdot \text{severity}(x) \;+\; w_2 \cdot \bigl(1 - \text{confidence}(n, x)\bigr) \;+\; w_3 \cdot \text{distance}\bigl(n, h(n)\bigr) \tag{PF}
\end{equation}
where $\text{severity}(x) \in [0, 1]$ is the assigned severity of exception $x$, $\text{distance}(n, h(n))$ is the number of hops from $n$ to the human anchor in the parent chain, and $w_1, w_2, w_3$ are provisioned constants satisfying $w_1 > w_2 > w_3 \geq 0$. The highest-priority trigger fires first.

\subparagraph{Secondary trigger: explicit \texttt{bailout\_signal}.} The Bounds Engine may emit an explicit \texttt{bailout\_signal} at any time, bypassing TC entirely:
\begin{equation}
\texttt{bailout\_signal}(x) \;\implies\; \text{T1 fires for node } n \text{ with exception } x. \tag{TC-secondary}
\end{equation}
This handles conditions — for example, a logical contradiction between two inputs within $R(n)$ — that the Bounds Engine recognizes as unresolvable even though both inputs are individually within the operating range.

% --------------------------------------------------------------
\subsection{Escalation Algorithm}\label{sec:escalation-algorithm}

When the state machine enters ESCALATING, the protocol executes the \texttt{Escalate} algorithm to propagate the exception to $h(n)$ along the parent chain. The algorithm is executed by the Bailout Monitor on the originating node; intermediate nodes execute a lightweight forwarding function with no autonomous evaluation.

\begin{algorithm}[htbp]
\caption{\texttt{Escalate($n$, $x$, $d$)}}
\label{alg:escalate}
\begin{algorithmic}[1]
\Procedure{Escalate}{$n, x, d$}
\State $\text{diag} \leftarrow \text{ComposeDiagnostic}(n, x, d)$
\State $\text{diag\_hash} \leftarrow \text{SHA256}(\text{diag})$
\State $\text{ledger\_entry} \leftarrow \text{CreateLedgerEntry}('ESCALATION\_START', n, x, \text{diag\_hash})$
\State $\text{FactLayer.Commit}(\text{ledger\_entry})$
\State $p \leftarrow \alpha(n)$
\State $T_{\text{prop}} \leftarrow T_{\text{init}}$
\State $\text{retries} \leftarrow 0$
\State $\text{pathway} \leftarrow \text{SelectPathway}(PHR)$
\State $\text{msg} \leftarrow \text{Serialize}(\text{diag}, \text{diag\_hash}, n)$
\State $\text{ack\_received} \leftarrow \texttt{false}$
\State
\While{$\neg\text{ack\_received} \land \text{retries} \leq N_{\max}$}
\State $\text{send\_result} \leftarrow \text{pathway.Send}(p, \text{msg}, T_{\text{prop}})$
\If{$\text{send\_result} = \texttt{ACK}$}
\State $\text{ack\_received} \leftarrow \texttt{true}$
\State $\text{RecordAck}(p, \text{diag\_hash})$
\State \textbf{goto} \texttt{WAIT\_DIRECTIVE}
\ElsIf{$\text{send\_result} = \texttt{TIMEOUT}$}
\State $\text{retries} \leftarrow \text{retries} + 1$
\State $T_{\text{prop}} \leftarrow \min(2 \cdot T_{\text{prop}},\ T_{\text{cap}})$
\State $\text{pathway} \leftarrow \text{SelectPathway}(PHR)$
\State $\text{RecordRetry}(p, \text{retries}, T_{\text{prop}})$
\EndIf
\EndWhile
\State
\If{$\neg\text{ack\_received} \land \text{retries} > N_{\max}$}
\State $\text{FireParentUnreachable}(n, x, \text{diag})$
\State $\text{TransitionState}(\text{ESCALATING}, \text{RESOLVED})$
\State \textbf{return}
\EndIf
\State
\State \texttt{WAIT\_DIRECTIVE:}
\State $\text{awaiting\_directive} \leftarrow \texttt{true}$
\While{$\text-awaiting\_directive$}
\State $\text{directive} \leftarrow \text{ReceiveDirective}(h(n), T_{\text{human}})$
\If{$\text{directive} = \texttt{ACK}$}
\State $\text{TransitionState}(\text{HUMAN\_ACKNOWLEDGED}, \text{IDLE})$
\State $\text-awaiting\_directive} \leftarrow \texttt{false}$
\ElsIf{$\text{directive} \in \{\texttt{RESOLVE}, \texttt{REJECT}\}$}
\State $\text{RecordDirective}(\text{directive}, h(n))$
\State $\text{TransitionState}(\text{HUMAN\_ACKNOWLEDGED}, \text{RESOLVED})$
\State $\text-awaiting\_directive} \leftarrow \texttt{false}$
\ElsIf{$\text{directive} = \texttt{TIMEOUT}$}
\State $\text{SendReminder}(h(n))$
\State $\text{TransitionState}(\text{HUMAN\_ACKNOWLEDGED}, \text{RESOLVED}, \texttt{unresolved})$
\State $\text-awaiting\_directive} \leftarrow \texttt{false}$
\EndIf
\EndWhile
\State \textbf{return}
\EndProcedure
\end{algorithmic}
\end{algorithm}

\subparagraph{\texttt{SelectPathway} function.} The \texttt{SelectPathway} function queries the Pathway Health Registry (PHR) and returns the functional pathway with the lowest estimated latency:
\begin{equation}
\text{SelectPathway}(PHR) \;=\; \underset{p \in \text{functional\_pathways}}{\arg\min}\ \text{latency\_estimate}(p) \tag{SP}
\end{equation}
Pathways marked \texttt{DEGRADED} after $K$ consecutive health-check failures are excluded from selection until a health check succeeds.

\subparagraph{\texttt{ComposeDiagnostic} function.} The diagnostic payload assembled at each escalation step includes the exception descriptor, the node's current operating range, confidence and severity assessments, the full parent chain history, a SHA-256 hash of the diagnostic content for tamper-evidence, and the \fact{} layer authorization state at time of escalation. This payload is appended, not overwritten, at each hop — preserving the complete forensic chain from originating node to human anchor.

% --------------------------------------------------------------
\subsection{Bypass Mechanism: Why the Bailout Protocol Bypasses All Machine Actors}\label{sec:bypass-mechanism}

The parent-chain bypass is the structural property that distinguishes the Bailout Protocol from conventional tiered oversight architectures. It ensures that when the state machine is in ESCALATING, the exception propagates to $h(n)$ without passing through any intermediate machine actor as a terminus. No machine actor in the parent chain may absorb, resolve, or suppress the exception.

\subparagraph{Formal bypass property.} The bypass is stated formally as:
\begin{equation}
\text{BYPASS}(n) \;\equiv\; \forall m \in \text{MachineActors} \setminus \{h(n)\}:\; m \notin \text{terminus}\bigl(\text{ESCALATING}(n)\bigr) \tag{BP}\label{eq:bypass-property}
\end{equation}
where $\text{terminus}(\text{ESCALATING}(n))$ denotes the set of actors that may be the final recipient of the ESCALATING state for node $n$. BYPASS requires that this set contains exactly one element: the human anchor $h(n)$. No machine actor may appear in this set under any execution condition.

\subparagraph{Why bypass is architecturally necessary.} The bypass is not an optimization — it is a structural necessity that follows directly from the upstream accountability guarantee of the \bxthree{} Framework. Any machine actor $m$ that could absorb an exception in ESCALATING state would become a terminal node for that exception, and the accountability chain would terminate at $m$ rather than at a human. If $m$ then fails, or if $m$'s resolution is incorrect or insufficient, no human principal is present in the chain to detect or correct the failure. The system compounds failure instead of surfacing it. This is the precise failure mode that the Bailout Protocol is designed to prevent.

Conventional escalation hierarchies — including the Tiered Agentic Oversight (TAO) model~\cite{kim2025tao} — permit this compounding by making human oversight a high-tier option at the top of a machine-tier hierarchy, where each intermediate tier may absorb the exception before passing it up. The human anchor in TAO is present but not architecturally required for every exception; absorption at an intermediate tier is permitted and invisible to external auditors. The Bailout Protocol forbids this absorption by architectural constraint: the human anchor is not one option among many — it is the \emph{only} permissible terminus for every unresolved exception. The parent-chain bypass skips intermediate nodes not solely for efficiency but because their presence in the terminus would break the accountability chain.

\subparagraph{Enforcement mechanisms.} The bypass is enforced by two cooperating mechanisms operating at different layers:

\begin{enumerate}
\item[\textbf{\fact{} Layer gate (enforcement at the authorization layer).}] The \fact{} layer is the only mechanism capable of authorizing state transitions in the \bxthree{} Framework. When a node $n$ is in BOUNDED state and produces a resolution $a$, the resolution is written to the authorization ledger only after \fact{} layer approval. The Bailout Monitor's pre-commit hook rejects any ledger entry for a resolution where the underlying trigger satisfies TC and $T_{\text{bounds}}$ has expired. This prevents the Bounds Engine from issuing a retroactive resolution after timeout to suppress warranted escalation. The \fact{} layer gate operates identically on all machine actors: there is no privileged machine actor with the ability to bypass this check.

\item[\textbf{Immutable parent-pointer and human-root architecture (enforcement at the data model layer).}] The parent pointer $\alpha(n)$ and the human-root identifier $h(n)$ are stored in \fact{} layer-managed memory and are read-only at the machine-actor level. A machine actor cannot redirect the parent chain, cannot insert itself as an intermediate resolver, and cannot modify the $humanRootId$ field that terminates the chain. Any attempted modification requires \purpose{} layer authorization with a corresponding ledger entry — which cannot be completed without human approval if the modification would affect the bailout terminus.
\end{enumerate}

Together, these two mechanisms ensure that BYPASS holds for all $n \in \text{nodes}$ under all execution conditions, including partial node failure, network partition, and deliberate attempt by a compromised machine actor to redirect the bailout pathway.

% --------------------------------------------------------------
\subsection{Timeout Handling}\label{sec:timeout-handling}

Timeouts in the Bailout Protocol serve a liveness function: they prevent the protocol from permanently stalling when a parent node is unreachable, a communication pathway is unresponsive, or the human anchor is temporarily unavailable. Each timeout has a defined scope, a defined scope-exit action, and an escalation path that prevents the protocol from deadlocking.

\subparagraph{Timeout definitions.}

\begin{enumerate}
\item[\textbf{$T_{\text{bounds}}$ — Bounds Engine resolution timeout.}] The maximum time a node may remain in BOUNDED state attempting autonomous resolution. Default: 60 seconds, provisioned per node class according to the expected latency of the node's resolution procedure. If $T_{\text{bounds}}$ expires without a \fact{} layer-approved resolution, the node transitions to BAIL\_PENDING $\rightarrow$ ESCALATING via T3 and T4. The \fact{} layer cannot extend $T_{\text{bounds}}$ unilaterally; any extension requires human authorization via a separate ledger entry recorded before $T_{\text{bounds}}$ expires.

\item[\textbf{$T_{\text{prop}}$ — Per-hop propagation timeout.}] The maximum time node $n$ waits for an acknowledgment from its parent $p$ after sending an exception via \texttt{SendToParent}. Default initial value: 5 seconds. If $T_{\text{prop}}$ expires without ACK, the algorithm retries with exponential backoff: $T_{\text{prop}} \leftarrow \min(2 \cdot T_{\text{prop}},\ T_{\text{cap}})$. This backoff continues until either an ACK is received or $N_{\max}$ retries have been exhausted at $T_{\text{cap}}$.

\item[\textbf{$T_{\text{cap}}$ — Maximum propagation timeout value.}] The ceiling value for the exponential backoff of $T_{\text{prop}}$. Default: 60 seconds. After $N_{\max}$ consecutive retries at $T_{\text{cap}}$ without ACK, the algorithm fires a \texttt{parent\_unreachable} timeout event, treated as a new trigger event propagating to $h(n)$ via the next available functional pathway tracked by the Pathway Health Registry.

\item[\textbf{$T_{\text{human}}$ — Human acknowledgment timeout.}] The maximum time the human anchor may remain in HUMAN\_ACKNOWLEDGED state without issuing a directive. Default: 300 seconds (5 minutes), configurable per deployment. If $T_{\text{human}}$ expires, the protocol issues a reminder notification to $h(n)$ and transitions the originating node to RESOLVED with an \texttt{unresolved} tag. The node becomes quiescent: no autonomous retry occurs, no further actions are attempted, and the \texttt{unresolved} tag signals that human re-engagement is required to clear the state. The protocol does not re-escalate automatically; automatic re-escalation would create a potential denial-of-service condition against the human anchor.
\end{enumerate}

\subparagraph{Pathway Health Registry (PHR).} The PHR is a runtime data structure maintained by the Bailout Monitor that tracks the availability, latency, and error rate of each provisioned communication pathway to the human anchor. Before each \texttt{SendToParent} call, \texttt{SelectPathway} queries the PHR and selects the pathway with the lowest estimated latency among currently functional pathways. Pathways that fail $K$ consecutive health checks (default $K = 3$) are marked \texttt{DEGRADED} and excluded from pathway selection until a health check succeeds.

The PHR is updated by a background health-ping thread that sends heartbeat messages along each provisioned pathway at intervals of $T_{\text{health}}$ (default: 30 seconds). The health-ping thread operates with the same scheduling priority as the Bailout Monitor and is not preemptible by agent tasks, ensuring that pathway status remains current even under heavy agent workload.

\subparagraph{Liveness guarantee.} Together, $T_{\text{bounds}}$, $T_{\text{prop}}$, $T_{\text{cap}}$, and $T_{\text{human}}$ define a finite upper bound on the wall-clock time from trigger firing (T1) to human acknowledgment under any combination of parent unavailability and pathway degradation:
\begin{equation}
T_{\max} \;=\; T_{\text{bounds}} \;+\; D \cdot \bigl(T_{\text{cap}} \cdot N_{\max} + T_{\text{health}}\bigr) \;+\; T_{\text{human}} \tag{T-max}
\end{equation}
where $D$ is the number of hops from the originating node to $h(n)$. This bound is a deployment parameter — not a protocol constant — derived from network topology and provisioned retry parameters. A deployment that cannot guarantee human acknowledgment within the required $T_{\max}$ under all failure scenarios must provision additional or higher-bandwidth pathways, or must reduce $T_{\text{prop}}$ and $T_{\text{cap}}$ to bring the worst-case bound within the required limit.
\end